\subsection{环的变换}

命题 7.--- 设 \(\rho: A \to B\) 是诺特局部环的局部同态，M 是有限生成 A-模，N 是有限生成 B-模。置 \(\overline{B} = B \otimes_A \kappa_A = B / \rho(m_A).B\)，以及 \(\overline{N} = N \otimes_B \overline{B} = N / \rho(m_A).N\)。我们有

\[
\dim_ {B} (M \otimes_ {A} N) \leqslant \dim_ {A} (M) + \dim_ {B} (\overline {{N}}),\tag{12}
\]

且当 N 在 A 上平坦时取等号。

不妨设 M 和 N 非零。根据 th. 1 的推论（n° 2），存在 \(m_A\) 的子集 S（分别为 \(m_B\) 的子集 T），使 \(\text{Card}(S) = \dim_A(M)\)（分别有 \(\text{Card}(T) = \dim_{\overline{B}}(\overline{N})\)），并且 M/SM 是有限长度的 A-模（分别地，\(\overline{N}/\overline{T}\overline{N}\) 是有限长度的 B-模）。有 \(\rho(S) \subset m_B\)。令 E 为 B-模 \(M \otimes_A N\)；B-模 \(E/(\rho(S).E + T.E)\) 同构于 M/SM \(\otimes_A N/TN\)，因而具有有限长度；根据 II, § 4, n° 4 的 prop. 18 和 19，有

\[
\begin{array}{r l} \operatorname{Supp} (\mathrm{M} / \mathrm{SM} \otimes_ {\mathrm{A}} \mathrm{N} / \mathrm{TN}) & = \operatorname{Supp} (\mathrm{N} / \mathrm{TN}) \cap {} ^ {a} \rho^ {- 1} (\operatorname{Supp} (\mathrm{M} / \mathrm{SM})) \\ & = \operatorname{Supp} (\mathrm{N} / \mathrm{TN}) \cap {} ^ {a} \rho^ {- 1} (\operatorname{Supp} (\kappa_ {\mathrm{A}})) \\ & = \operatorname{Supp} (\mathrm{N} / \mathrm{TN} \otimes_ {\mathrm{A}} \kappa_ {\mathrm{A}}) = \operatorname{Supp} (\overline {{\mathrm{N}}} / \mathrm{T} \overline {{\mathrm{N}}}) = \left\{\mathfrak {m} _ {\mathrm{B}} \right\}. \end{array}
\]

根据 th. 1，得到不等式

\[
\dim_ {\mathrm{B}} (\mathrm{E}) \leqslant \operatorname{Card} (\mathrm{S}) + \operatorname{Card} (\mathrm{T}) = \dim_ {\mathrm{A}} (\mathrm{M}) + \dim_ {\overline {{\mathrm{B}}}} (\overline {{\mathrm{N}}}).
\]

现在假设 N 在 A 上平坦，并证明公式 (12) 中取等号。设 a（分别为 b）是 M（分别为 N）的湮灭子。根据 II, § 4, n° 4 的 prop. 18 和 19，有

\[
\begin{array}{r l} \operatorname{Supp} (E) & = \operatorname{Supp} (N) \cap {} ^ {a} \rho^ {- 1} (\operatorname{Supp} (M)) \\ & = V (b) \cap {} ^ {a} \rho^ {- 1} (V (a)) = V (b + \rho (a). B). \end{array}
\]

因此有

\[
\dim_ {B} (M \otimes_ {A} N) = \dim (B / (b + \rho (a). B))\tag{13}
\]

并且

\[
\dim_ {\mathbf {A}} (\mathbf {M}) + \dim_ {\overline {{\mathbf {B}}}} (\overline {{\mathbf {N}}}) = \dim (\mathrm{A} / \mathfrak {a}) + \dim (\mathbf {B} / (\mathfrak {b} + \rho (m _ {\mathbf {A}}). \mathbf {B}))  .\tag{14}
\]

置 A' = A/a，B' = B/(b + ρ(a).B)，并令 ρ':A' → B' 为由 ρ 经过取商得到的局部同态。由于 N 的湮灭子为 b，N 是一个支集等于 Spec(B/b) 的有限生成 B/b-模，并且在 A 上平坦。因此，由 ρ 导出的同态 A → B/b 具有 § 2, n° 1 的性质 (PM)（loc. cit., remark 3）；由标量扩张（loc. cit., prop. 1, a)），可推出 ρ' 具有性质 (PM)。根据 loc. cit. 的 prop. 2，有

\[
\dim (\mathbf {B} ^ {\prime}) \geqslant \dim (\mathbf {A} ^ {\prime}) + \dim (\mathbf {B} ^ {\prime} / \rho^ {\prime} (m _ {\mathbf {A} ^ {\prime}}). \mathbf {B} ^ {\prime})\tag{15}
\]

又由于环 \(\mathbf{B}'/\rho'(\mathfrak{m}_{\mathbf{A}'})\mathbf{B}'\) 同构于 \(\mathbf{B}/(\mathfrak{b} + \rho(\mathfrak{m}_{\mathbf{A}})\mathbf{B})\)，我们的断言由公式 (13)、(14) 和 (15) 得出。

推论 1. --- 设 \(\rho: A \to B\) 是诺特局部环的局部同态。

\begin{enumerate}
\def\labelenumi{\alph{enumi})}
\tightlist
\item
  我们有
\end{enumerate}

\[
\dim (\mathbf {B}) \leqslant \dim (\mathbf {A}) + \dim (\mathbf {B} \otimes_ {\mathbf {A}} \kappa_ {\mathbf {A}})\tag{16}
\]

且当 ρ 使 B 成为平坦的 A-模时取等号。

\begin{enumerate}
\def\labelenumi{\alph{enumi})}
\setcounter{enumi}{1}
\tightlist
\item
  假设 \(\rho\) 使 B 成为平坦的 A-模，并设 S 是 \(\mathfrak{m}_{\mathbf{A}}\) 的子集。则 S 是 A 的截切集，当且仅当 \(\rho(\mathbf{S})\) 是 B 的截切集。
\end{enumerate}

断言 \(a)\) 是 prop. 7 在 \(\mathbf{M} = \mathbf{A}\)、\(\mathbf{N} = \mathbf{B}\) 时的特殊情形。

证明 b)。在上述假设下，有

\[
\dim (\mathbf {B}) = \dim (\mathbf {A}) + \dim (\overline {{\mathbf {B}}})\tag{17}
\]

其中 \(\overline{\mathbf{B}} = \mathbf{B} \otimes_{\mathbf{A}} \kappa_{\mathbf{A}} = \mathrm{B} / \rho(\mathfrak{m}_{\mathbf{A}}) . \mathbf{B}\)。由于 \(\rho\) 是单射（I, § 3, no. 5, prop. 9），有

\[
\operatorname{Card} (\rho (S)) = \operatorname{Card} (S).\tag{18}
\]

最后，B′ = B/ρ(S). B 是 A′ = A/SA 上的平坦模，因而

\[
\dim (\mathrm{B} / \rho (\mathrm{S}). \mathrm{B}) = \dim (\mathrm{A} / \mathrm{SA}) + \dim (\overline {{\mathrm{B}}})\tag{19}
\]

因为 \(\mathbf{B}^{\prime}/\rho(\mathfrak{m}_{\mathbf{A}^{\prime}})\mathbf{B}^{\prime}\) 同构于 \(\overline{B}\)。我们的断言立即由公式 (17)、(18) 和 (19) 得到。

推论 2. — 设 ρ : A → B 是诺特环的同态。则

\[
\dim (\mathbf {B}) \leqslant \dim (\mathbf {A}) + \sup _ {\mathfrak {p} \in \operatorname{Spec} (\mathbf {A})} \dim (\mathbf {B} \otimes_ {\mathbf {A}} \kappa (\mathfrak {p})).
\]

设 q 是 B 的素理想，且 \(\mathfrak{p} = \rho^{-1}(\mathfrak{q})\)。由 cor. 1，有 \(\dim(B_{\mathfrak{q}}) \leqslant \dim(A_{\mathfrak{p}}) + \dim(B_{\mathfrak{q}} \otimes_{A} \kappa(\mathfrak{p}))\)。由此根据 § 1, no. 3 的 prop. 6, b) 得到不等式 \(\dim(B_{\mathfrak{q}}) \leqslant \dim(A) + \dim(B \otimes_{A} \kappa(\mathfrak{p}))\)，再由 § 1, no. 3 的 prop. 8 得到结论。

推论 3. --- 设 A 为诺特环，n 为整数 ≥ 0。则

\[
\dim (\mathrm{A} [ \mathrm{X} _ {1},..., \mathrm{X} _ {n} ]) = \dim (\mathrm{A}) + n.
\]

置 \(\mathbf{B} = \mathbf{A}[\mathbf{X}_1, \ldots, \mathbf{X}_n]\)。对于每个素理想 \(\mathfrak{p}\)，即 \(\mathbf{A}\) 的素理想，环 \(\mathbf{B} \otimes_{\mathbf{A}} \kappa(\mathfrak{p})\) 是域上的 \(n\) 元多项式环，因而维数为 \(n\)（\(\S 2\)，no 4, cor. 1 to th. 3）。由 cor. 2，有 \(\dim(\mathbf{B}) \leqslant \dim(\mathbf{A}) + n\)；反向不等式由 \(\S 1\) 的 Example 4, no. 3 得到。

推论 4. --- 设 \(\rho: A \to B\) 是诺特环的同态，a 是 A 的理想。我们有不等式

\[
\operatorname{ht} (\rho (\mathfrak {a}). \mathbf {B}) \leqslant \operatorname{ht} (\mathfrak {a})\tag{20}
\]

若映射 \(^a\rho : \text{Spec}(\mathbf{B}) \to \text{Spec}(\mathbf{A})\) 是满射。若 B 是忠实平坦的 A-模，则有 \(\operatorname{ht}(\rho(\mathfrak{a}).\mathbf{B}) = \operatorname{ht}(\mathfrak{a})\)。

若 B 在 A 上忠实平坦，则 \(^{a}\rho\) 是满射（II, § 2, n° 5, cor. 4 to prop. 11），并且有 ht(a) ≤ ht(ρ(a).B)（§ 2, n° 1, corollary of prop. 2）。

因此，只需在 \(^a\rho\) 为满射的假设下证明不等式 (20)。设 \(\mathfrak{p}\) 是 A 的素理想，满足 \(\mathfrak{a} \subset \mathfrak{p}\) 且 \(\mathrm{ht}(\mathfrak{a}) = \mathrm{ht}(\mathfrak{p})\)（\(\S 1, n^o 3, \text{prop. } 7\)）。令 \(\overline{\mathbf{B}} = \mathbf{B} \otimes_{\mathbf{A}} \kappa(\mathfrak{p})\)，并以 \(h\) 表示从 B 到 \(\overline{\mathbf{B}}\) 的典范同态。若 \(X = ^a\rho^{-1}(\mathfrak{p})\) 是位于 \(\mathfrak{p}\) 之上的 B 的素理想组成的非空集合，我们知道（\(\S 2, n^o 1, \text{lemme 1}\)），映射 \(^ah\) 是从 \(\operatorname{Spec}(\overline{\mathbf{B}})\) 到 \(\operatorname{Spec}(\mathbf{B})\) 的子空间 X 的同胚。设 \(\mathfrak{q}\) 是 \(^ah\) 下来自 \(\overline{\mathbf{B}}\) 的一个极小素理想的像；有 \(\dim(\mathbf{B}_{\mathfrak{q}} \otimes_{\mathbf{A}} \kappa(\mathfrak{p})) = 0\)，并且 Cor. 1 蕴含不等式 \(\dim(\mathbf{B}_{\mathfrak{q}}) \leqslant \dim(\mathbf{A}_{\mathfrak{p}})\)。最后，

\[
\operatorname{ht} (\rho (a). B) \leqslant \operatorname{ht} (q) = \dim (B _ {q}) \leqslant \dim (A _ {p}) = \operatorname{ht} (p) = \operatorname{ht} (a),
\]

从而得到该推论。

命题 8.--- 设 A 为诺特环，a 为 A 的理想，M 为有限生成 A-模，\(\hat{\mathbf{A}}\) 和 \(\hat{\mathbf{M}}\) 分别是 A 和 M 关于 a-进拓扑的分离完备化。

\begin{enumerate}
\def\labelenumi{\alph{enumi})}
\item
  设 m 是包含 a 的 A 的素理想，且 \(\hat{\mathfrak{m}} = \mathfrak{m}\hat{\mathbf{A}}\)。则 \(\hat{\mathfrak{m}}\) 是 \(\hat{\mathbf{A}}\) 的素理想，并且有 \(\dim_{\hat{\mathbf{A}}_{\hat{\mathfrak{m}}}}(\hat{\mathbf{M}}_{\hat{\mathfrak{m}}}) = \dim_{\mathbf{A}_{\mathfrak{m}}}(\mathbf{M}_{\mathfrak{m}})\)。
\item
  我们有 \(\dim_{\hat{\mathbf{A}}}(\hat{\mathbf{M}}) = \sup_{\mathfrak{m}} \dim_{\mathbf{A}_{\mathfrak{m}}}(\mathbf{M}_{\mathfrak{m}})\)，其中 \(\mathfrak{m}\) 遍历包含 a 的 A 的素理想
\end{enumerate}

（分别地，极大理想）组成的集合。特别地，有 \(\dim_{\hat{\mathbf{A}}}(\hat{\mathbf{M}}) \leqslant \dim_{\mathbf{A}}(\mathbf{M})\)。

\begin{enumerate}
\def\labelenumi{\alph{enumi})}
\item
  由于 \(\hat{\mathbf{A}}/\hat{\mathfrak{m}}\) 可识别为 \(\mathbf{A}/\mathfrak{m}\)，\(\hat{\mathfrak{m}}\) 是 \(\hat{\mathbf{A}}\) 的素理想。由 III, § 3, n° 4 的 th. 3，\(\hat{\mathbf{A}}\) 在 \(\mathbf{A}\) 上平坦，因而 \(\hat{\mathbf{A}}_{\hat{\mathfrak{m}}}\) 在 \(\mathbf{A}_{\mathfrak{m}}\) 上平坦。此外，\(\mathbf{A}\) 到 \(\hat{\mathbf{A}}\) 的典范映射诱导出 \(\mathbf{A}/\mathfrak{a}\) 到 \(\hat{\mathbf{A}}/\mathfrak{a}\hat{\mathbf{A}}\) 的同构，因而也诱导出 \(\mathbf{A}_{\mathfrak{m}}/\mathfrak{mA}_{\mathfrak{m}}\) 到 \(\hat{\mathbf{A}}_{\hat{\mathfrak{m}}}/\mathfrak{m}\hat{\mathbf{A}}_{\hat{\mathfrak{m}}}\) 的同构。将 prop. 7 应用于环 \(\mathbf{A}_{\mathfrak{m}}\) 和 \(\hat{\mathbf{A}}_{\hat{\mathfrak{m}}}\)，以及模 \(\mathbf{M}_{\mathfrak{m}}\) 和 \(\hat{\mathbf{A}}_{\hat{\mathfrak{m}}}\)，即可得到结论；模 \(\mathbf{M}_{\mathfrak{m}} \otimes_{\mathbf{A}_{\mathfrak{m}}} \hat{\mathbf{A}}_{\hat{\mathfrak{m}}}\) 同构于 \(\hat{\mathbf{M}}_{\hat{\mathfrak{m}}}\)（III, loc. cit. and prop. 8）。
\item
  由 III, § 3, n° 4 的 prop. 8，映射 m ↦ m̂ 是从包含 a 的 A 的极大理想集合到 Â 的极大理想集合的双射。断言 b) 由此以及 § 1, n° 4 的 prop. 9 得到。
\end{enumerate}

推论 1. --- 设 A 是 Zariski 环（III, § 3, n° 3, def. 2）。对于每个有限生成 A-模 M，有 \(\dim_{\hat{\mathbf{A}}}(\hat{\mathbf{M}}) = \dim_{\mathbf{A}}(\mathbf{M})\)。

事实上，A 的拓扑是 a-进拓扑，其中 a 是包含于 A 的根基中的理想（loc. cit.），也就是说，a 包含于 A 的每个极大理想中。因此只需应用 prop. 8 的断言 b)。

推论 2. --- 设 A 为诺特环，a 为 A 的理想，\(\hat{A}\) 为 A 关于 a-进拓扑的分离完备化。我们有 dim( \(\hat{A}\) ) \(\leqslant\) dim(A)；当 A 是局部环且 \(a\) 不同于 A 时取等号。

推论 3. --- 设 A 为诺特环，n 为整数 ≥ 0。则

\[
\dim (\mathrm{A} [ [ \mathrm{X} _ {1},..., \mathrm{X} _ {n} ] ]) = \dim (\mathrm{A}) + n.
\]

环 A{[}{[}X₁, \ldots, Xₙ{]}{]} 是多项式环 A{[}X₁, \ldots, Xₙ{]} 关于 \(\mathfrak{a}\)-进拓扑的分离完备化，其中 \(\mathfrak{a}\) 是由 X₁, \ldots, Xₙ 生成的理想；因而有

\[
\dim (\mathrm{A} [ [ \mathrm{X} _ {1},..., \mathrm{X} _ {n} ] ]) \leqslant \dim (\mathrm{A} [ \mathrm{X} _ {1},..., \mathrm{X} _ {n} ])
\]

根据 cor. 2。我们还有

\[
\dim (\mathrm{A} [ \mathrm{X} _ {1},..., \mathrm{X} _ {n} ]) = \dim (\mathrm{A}) + n
\]

根据 prop. 7 的 cor. 3。最后，有

\[
\dim (\mathrm{A}) + n \leqslant \dim (\mathrm{A} [ [ \mathrm{X} _ {1},..., \mathrm{X} _ {n} ] ])
\]

根据 § 1, n° 3 的 example 4。

备注。--- 1) 设 A 为诺特环，a 为 A 的理想。假设 A 关于 a-进拓扑是分离且完备的，并考虑受限形式级数环 R = A \{ X \(_{1}\) , \ldots, X \(_{n}\) \}（III, § 4, n° 2, def. 2）。我们有 dim(R) = dim(A) + n。事实上，R 是环 B = A{[}X \(_{1}\) , \ldots, X \(_{n}\) {]} 关于 aB-进拓扑的完备化，因而 dim(R) ≤ dim(A{[}X \(_{1}\) , \ldots, X \(_{n}\) {]}) = dim(A) + n。反向不等式的证明与形式级数的情形相同。

\begin{enumerate}
\def\labelenumi{\arabic{enumi})}
\setcounter{enumi}{1}
\tightlist
\item
  设 A 是诺特局部环，并将其识别为完备化 \(\hat{\mathbf{A}}\) 的子环；设 B 是包含 A 的 \(\hat{\mathbf{A}}\) 的子环。假设 B 是诺特局部环，且有 \(m_{\mathbf{A}}B = m_{\mathbf{B}}\)。则 A 到 B 的典范嵌入延拓为 \(\hat{\mathbf{A}}\) 到 B 的完备化 \(\hat{\mathbf{B}}\) 的同构（III, § 3, no. 5, prop. 11），因而 dim(B) = dim(A)（prop. 8 的 cor. 2）。特别地，这适用于下述情形。* 设 k 是完备的非离散赋值域，A 是多项式环 \(k[\mathrm{X}_1, ..., \mathrm{X}_n]\) 在由 \(\mathrm{X}_1, ..., \mathrm{X}_n\) 生成的素理想处的局部环，B 是系数在 k 中的 \(\mathrm{X}_1, ..., \mathrm{X}_n\) 的收敛幂级数环。则上述假设得到满足，从而有 dim(B) = n。*
\end{enumerate}

